\documentclass[11pt]{article}

\usepackage[margin=1in]{geometry}
\usepackage{amsmath, amssymb, amsthm}
\usepackage{mathtools}
\usepackage{bm}
\usepackage{booktabs}

\DeclareMathOperator*{\clip}{clip}
\DeclareMathOperator{\sign}{sgn}
\newcommand{\Ehat}{\widehat{\mathbb{E}}}
\newcommand{\dif}{\mathrm{d}}

\title{Sign-Handling in $\Delta_z^{(W)} = \sign(s_z)\cdot\clip(\min(s_z^+,s_z^-),\delta_z)$ \\
\large Trade-off Analysis of Three (Four) Variants}
\author{}
\date{}

\begin{document}
\maketitle

\begin{abstract}
The BSPO Wasserstein clipping operator
$\Delta_z^{(W)} = \sign(s_z)\cdot\clip(\min(s_z^+,s_z^-),\delta_z)$ carries
a non-differentiable $\sign$ factor. We enumerate three candidate
treatments (with two sub-interpretations of the third) and compare them
on: mathematical faithfulness, gradient behaviour at the on-policy fixed
point, gradient behaviour just off-policy, hyperparameter cost, and
expected effect on final optimisation. The key observation is that the
magnitude factor $\clip(\min(s_z^+,s_z^-),\delta_z)$ is identically zero
at the on-policy state, \emph{pre-annihilating} any sign contribution
regardless of how the sign is implemented. We conclude that
$\verb|torch.sign(s).detach()|$ is the unique choice consistent with the
mathematical derivative ($\dif\sign/\dif s = 0$ almost everywhere) and
recommend it as the default, with a concrete empirical protocol to
validate that the sign path is informationally inert.
\end{abstract}

% =============================================================================
\section{The three variants defined precisely}
\label{sec:variants}

We need $\sign(s_z)$ in forward; the question is what
$\dif\sign/\dif s$ returns to autograd in backward. Three variants are
considered.

\subsection*{V1 --- Stop-loss / \texttt{.detach()}}

\[
\sign_1(s) \;=\; \sign(s),\qquad
\frac{\dif\sign_1}{\dif s}(s) \;\equiv\; 0 \text{ (autograd override)}.
\]

Torch realisation: \verb|sign_s = torch.sign(s).detach()|.
Backward through the sign is severed; gradient through
$\Delta_z^{(W)}$ comes only from the magnitude factor
$\clip(\min(s_z^+, s_z^-), \delta_z)$.

\subsection*{V2 --- Continuous sign surrogate}

Two common smooth replacements:
\[
\sign_{2a}(s) \;=\; \frac{s}{\sqrt{s^2+\varepsilon^2}},
\qquad
\sign_{2b}(s) \;=\; \tanh(\beta s).
\]

For the $\varepsilon$-smoothed form,
\[
\frac{\dif\sign_{2a}}{\dif s}(s) = \frac{\varepsilon^2}{(s^2+\varepsilon^2)^{3/2}},
\qquad
\left.\frac{\dif\sign_{2a}}{\dif s}\right|_{s=0} = \frac{1}{\varepsilon}.
\]

Torch: \verb|sign_s = s / torch.sqrt(s*s + eps*eps)|. Fully
differentiable, no $\verb|.detach()|$.

\subsection*{V3 --- ``gradient $=1$ at $0$, $=0$ elsewhere''}

This prompt admits multiple interpretations.

\textbf{V3a --- literal single-point injection.} Define a custom autograd
node with forward $\sign(s)$ and backward $\dif\sign/\dif s = 1$ iff
$s = 0$, else $0$. In \texttt{fp32}/\texttt{fp64}, $s=0$ is measure-zero
and in practice almost never hit; this variant is \emph{numerically
indistinguishable from V1}.

\textbf{V3b --- local straight-through estimator (local STE).}
Inject gradient in an $\varepsilon$-neighborhood of $0$:
$\dif\sign/\dif s = 1$ for $|s| < \varepsilon$, $0$ otherwise. Torch:
\begin{verbatim}
near    = (s.abs() < eps).float()
sign_s  = torch.sign(s).detach() + near * (s - s.detach())
# forward = sgn(s); backward = 1 inside |s|<eps, 0 outside
\end{verbatim}

\textbf{V3c --- full STE (identity backward).}
Forward is $\sign(s)$, backward is the identity:
$\dif\sign/\dif s \equiv 1$. Torch: \verb|sign_s = torch.sign(s) + s - s.detach()|.
Common in discrete-variable / quantised NN literature.

% =============================================================================
\section{On-policy fixed-point analysis}
\label{sec:onpolicy}

At the on-policy state, $r_t = 1$ for all $t \in \tau$, hence
$s_\tau^+ = s_\tau^- = 0$, $s_\tau = 0$,
$\min(s_\tau^+, s_\tau^-) = 0$, $\clip(\min,\delta_\tau)=0$, and
$\Delta_\tau^{(W)} = 0$ regardless of the sign convention. The objective
contribution vanishes trivially.

The question is the \emph{gradient} at this state. Apply the product rule:
\begin{align}
\frac{\partial (\Delta_z^{(W)} \hat A_z)}{\partial \theta}
 &= \hat A_z\cdot\Bigl[
 \sign(s_z)\,\frac{\partial\,\clip(\min(s_z^+,s_z^-),\delta_z)}{\partial \theta}
 \nonumber \\
 &\qquad + \clip(\min(s_z^+,s_z^-),\delta_z)\,\frac{\partial \sign(s_z)}{\partial \theta}\Bigr].
 \label{eq:product-rule}
\end{align}

At on-policy:
\begin{itemize}
\item $\sign(0) = 0$ (by convention),
\item $\clip(\min,\delta_z)\big|_{s^{\pm}=0} = 0$,
\item $\partial\min/\partial\theta\big|_{s^{\pm}=0} = 0$ (any subgradient
works; the minimum at the boundary is $0$),
\item $\partial s/\partial \theta \neq 0$ in general (this is
$\partial(\log\pi_\theta - \log\pi_{\theta_{\text{old}}})/\partial\theta$).
\end{itemize}

Substituting into \eqref{eq:product-rule}:
\[
\left.\frac{\partial (\Delta_z^{(W)} \hat A_z)}{\partial \theta}\right|_{\text{on-policy}}
= \hat A_z\cdot\bigl[\,0\cdot 0 + 0\cdot \tfrac{\partial \sign_v}{\partial\theta}\,\bigr]
= 0.
\]

\begin{theorem}[Pre-annihilation at on-policy]
\label{thm:preannihilation}
For every sign-handling variant $v \in \{1, 2a, 2b, 3a, 3b, 3c\}$, and any
finite-valued $\partial\sign_v/\partial\theta$, the gradient of
$\Delta_z^{(W)}\hat A_z$ w.r.t.\ policy parameters $\theta$ is
identically zero at the exact on-policy state.
\end{theorem}

\begin{proof}
From \eqref{eq:product-rule}, both product terms carry a factor that
vanishes on-policy: the first is multiplied by $\sign(0)=0$, the second
by $\clip(\min,\delta_z)=0$. Any finite backward of
$\partial\sign_v/\partial\theta$ is multiplied by the zero magnitude.
\end{proof}

\textbf{Interpretation.} The sign path is \emph{informationally redundant}
at on-policy: whatever gradient the sign backward might wish to contribute
is multiplied by a magnitude factor that the definition of $\Delta^{(W)}$
already forces to zero. All six sign conventions produce identical
behaviour at the on-policy fixed point. Differences among them arise only
in the open neighbourhood $0 < |s| \leq \min(s^+, s^-)$.

% =============================================================================
\section{Off-policy regime and gradient-path decomposition}
\label{sec:offpolicy}

Set $m := \clip(\min(s_z^+, s_z^-), \delta_z) > 0$. Decompose
\eqref{eq:product-rule} into a \emph{magnitude path} (first term) and a
\emph{sign path} (second term):
\[
\frac{\partial (\Delta_z^{(W)} \hat A_z)}{\partial \theta}
= \underbrace{\sign(s_z)\,\hat A_z\cdot\frac{\partial m}{\partial \theta}}_{\text{magnitude path}}
+ \underbrace{m\,\hat A_z\cdot\frac{\partial \sign_v(s_z)}{\partial \theta}}_{\text{sign path}}.
\]

The magnitude path is identical across all six variants (they differ only
in the $\sign_v$ sub-expression, not in $m$). The sign path is what
separates them.

\subsection{Sign-path gradient at $s_z \to 0^+$ (just off-policy)}

$\partial\sign_v/\partial\theta = (\dif\sign_v/\dif s)\cdot\partial s/\partial\theta$.
Using the derivatives recorded in \S\ref{sec:variants}:

\begin{center}
\renewcommand{\arraystretch}{1.3}
\begin{tabular}{l l l}
\toprule
variant & $\dif\sign_v/\dif s$ at $s\to 0^+$ & sign-path gradient magnitude \\
\midrule
V1 (detach) & $0$ & $0$ \\
V2a (smooth, $s/\sqrt{s^2+\varepsilon^2}$) & $1/\varepsilon$ & $m\hat A_z / \varepsilon \cdot |\partial s/\partial\theta|$ \\
V2b (smooth, $\tanh(\beta s)$) & $\beta$ & $\beta m\hat A_z\cdot|\partial s/\partial\theta|$ \\
V3a (grad $=1$ at $s=0$ only) & $0$ a.e.\ (meas.~zero support) & $0$ in floating-point \\
V3b (local STE, $|s|<\varepsilon$) & $1$ & $m\hat A_z\cdot|\partial s/\partial\theta|$ \\
V3c (full STE) & $1$ & $m\hat A_z\cdot|\partial s/\partial\theta|$ \\
\bottomrule
\end{tabular}
\end{center}

\textbf{Consequences.}
V1 and V3a contribute zero to the sign path, matching the mathematical
a.e.\ derivative of $\sign$. V2a can produce a \emph{spike} of gradient of
size $m\hat A_z/\varepsilon$ --- unbounded as $\varepsilon\to 0$ --- near
the on-policy manifold. V2b is bounded by $\beta m$ but still biased. V3b
and V3c inject a constant gradient of $m\hat A_z\cdot|\partial s/\partial\theta|$
through the sign path regardless of whether the mathematical derivative
permits it.

% =============================================================================
\section{Trade-off matrix}
\label{sec:matrix}

Consolidating \S\ref{sec:variants}--\S\ref{sec:offpolicy}:

\begin{center}
\renewcommand{\arraystretch}{1.3}
\resizebox{\textwidth}{!}{%
\begin{tabular}{p{4.5cm} l l l l l}
\toprule
dimension & V1 detach & V2 smooth & V3a literal & V3b local STE & V3c full STE \\
\midrule
math faithfulness to $\dif\sign/\dif s = 0$ a.e. & exact & biased $\neq 0$ near $s=0$ & exact & biased inside $|s|<\varepsilon$ & heavily biased \\
sign-path gradient near $s=0$ & $0$ & $\sim m\hat A/\varepsilon$ & $0$ & $m\hat A\cdot\partial s/\partial\theta$ & $m\hat A\cdot\partial s/\partial\theta$ \\
magnitude-path gradient & $\sign(s)\hat A\cdot\partial m/\partial\theta$ & same (up to smoothed sign) & same & same & same \\
signal at ``just off-policy'' & small (mag is small too) & potentially large ($1/\varepsilon$ spike) & small & medium (bump active) & large (always pushes $s$) \\
hyperparameter & none & $\varepsilon$ or $\beta$ & none & $\varepsilon$ & none \\
explosive grad risk & none & high for small $\varepsilon$ & none & bounded by $\varepsilon$ & bounded, constant \\
biases policy away from $s\!=\!0$? & no & yes (via $m/\varepsilon$) & no & yes (inside bump) & yes (pushes $s$ up) \\
\texttt{.detach()} required? & yes (definitional) & no & yes on outer path & yes on outer & yes on outer \\
extra fwd-pass ops & $0$ & $\sim 2$ (sqrt, div) & $0$ & mask op & tiny (add/sub) \\
unbiased estimator of $\dif\sign/\dif s$? & yes (zero is truth) & no & yes & no & no \\
\bottomrule
\end{tabular}%
}
\end{center}

Three take-aways:

\begin{enumerate}
\item V1 is the only variant whose gradient matches the mathematical
a.e.\ derivative of $\sign$. Every other variant \emph{adds} a
non-mathematical gradient signal through the sign path, which is a bias.
\item V2a's bias is potentially catastrophic for small $\varepsilon$:
the $1/\varepsilon$ gradient at $s=0$, multiplied by a nonzero
$m\hat A$, injects unbounded gradient spikes near the on-policy
manifold. Avoiding this would require $\varepsilon \sim \delta_z$, at
which point the surrogate barely approximates $\sign$.
\item V3a is V1 in disguise. V3b and V3c are biased-on-purpose variants
that are appropriate only in settings where a deliberate push off the
on-policy manifold is desired.
\end{enumerate}

% =============================================================================
\section{Why V1 is the correct default --- the on-policy argument, revisited}
\label{sec:whyV1}

Theorem~\ref{thm:preannihilation} says: \emph{the magnitude factor
annihilates any sign contribution at on-policy}. Two corollaries follow.

\begin{corollary}[Sign-path informational redundancy]
Near the on-policy manifold where $\|s\|_\infty \to 0$,
$m \to 0$ at the same rate. Any sign-path gradient is thus
multiplied by a quantity that is itself going to zero. A well-scaled
training run does not need the sign path to contribute gradient near
on-policy; the magnitude path already handles the vanishing correctly.
\end{corollary}

\begin{corollary}[Non-mathematical sign-gradient is a bias, not a feature]
What V2/V3b/V3c's sign-path gradient actually does is push the policy
\emph{away} from the on-policy manifold. This is:
(i) not required by the math (Theorem~\ref{thm:preannihilation} shows the
objective doesn't need this push);
(ii) inconsistent with standard PG/PPO, which treat the sampling state
as the desired state to return to;
(iii) a modification of the objective, not a faithful implementation of it.
\end{corollary}

Thus V1 (detach) is the unique variant that (a) preserves the
$\Delta^{(W)} = \sign\cdot m$ factorisation as a mathematical statement
and (b) makes the backward chain consistent with the a.e.\ derivative of
$\sign$. The claim ``\texttt{.detach()} is a trick'' is incorrect: it is
the \emph{declaration} that the sign is a categorical direction selector,
and in doing so it realises the mathematically correct zero derivative
exactly, avoiding the NaN/undefined that a naive autograd might emit at
$s=0$.

% =============================================================================
\section{Expected effect on final optimisation performance}
\label{sec:effect}

How large is the effect of choosing V1 vs V2 vs V3? Theory alone gives us
\emph{bounds} on the effect, not its magnitude.

\subsection{Upper bound on the sign-path contribution}

For any step where $|s_z| > \varepsilon$ (outside V2/V3b's bump/spike
region), all variants give approximately identical gradients, because
either $\dif\sign_v/\dif s$ has decayed to zero (V2, V3b) or it was zero
to begin with (V1, V3a).

For steps with $|s_z| < \varepsilon$:
\[
\bigl\|\text{sign-path gradient}\bigr\|_{\infty}
\;\leq\; \frac{1}{\varepsilon}\cdot m\cdot|\hat A_z|\cdot\left\|\frac{\partial s}{\partial\theta}\right\|_\infty
\;\leq\; \frac{\delta_z\cdot\max|\hat A|\cdot\|\partial s/\partial\theta\|_\infty}{\varepsilon}.
\]

For the reference Phase-1 ablation ($\delta_z = 3\times 10^{-4}$,
$\max|\hat A|\approx 3$, $\|\partial s/\partial\theta\|_\infty$ bounded
by the model's per-token log-prob sensitivity), the worst-case
sign-path gradient under V2a with $\varepsilon = 10^{-4}$ is of the same
order as $|\partial s/\partial\theta|\cdot\delta_z/\varepsilon \approx 3$,
i.e.\ it \emph{is} a dominant contribution when it fires. But it fires
only in a vanishing-measure subset of steps (those with $|s_\tau| <
\varepsilon$), so its average effect on training is bounded by that
fraction of steps.

\subsection{Expected magnitude, heuristic}

In practice:
\begin{itemize}
\item After the first policy update, the ratio distribution spreads and
$\min(s_\tau^+, s_\tau^-)$ is almost always $> \varepsilon$ for sensible
choices (e.g.\ $\varepsilon < 10^{-3}$). The sign path is inactive for
V2, V3b most of the time.
\item When $|s_\tau|$ does get small, it is often because the
trajectory's ratios are symmetric around $1$, in which case
$\min(s_\tau^+, s_\tau^-)$ is also small but not zero --- so the
magnitude path gives a nonzero gradient anyway.
\item Empirically, Phase-1 v2 runs under V1 show
$\texttt{actor/bspo\_s\_tau\_abs\_mean}$ rising from $5.7\!\times\!10^{-4}$
at step $1$ to $\sim 0.2$ by step $60$ (see
\texttt{dev\_notes\_ablation.md}). The magnitude path alone is driving
the off-policy excursion; no sign-surrogate help is needed.
\end{itemize}

Expected effect on final validation accuracy: \textbf{$<1\%$ difference}
between V1 and V2 on mathematical reasoning benchmarks at standard BSPO
hyperparameters, with V1 more stable (no gradient spikes).

% =============================================================================
\section{Empirical measurement protocol}
\label{sec:measure}

We rank measurement strategies from cheapest to most rigorous.

\subsection*{(M1) Probe metrics during training (zero extra compute)}

Add two scalars to the training log:
\begin{itemize}
\item $\|\text{grad}_{\text{sgn-path}}\|$: norm of the sign-path
contribution, computed via a one-line autograd hook that captures
$\partial(m\hat A\cdot\sign_v(s))/\partial\theta$ specifically.
\item $\|\text{grad}_{\text{mag-path}}\|$: norm of the magnitude-path
contribution, captured symmetrically.
\end{itemize}

Diagnostic: if $\|\text{grad}_{\text{sgn}}\| \ll \|\text{grad}_{\text{mag}}\|$
throughout training, the sign path is informationally inert and all three
variants are indistinguishable. Under V1 the sign-path norm is exactly
$0$; this is a sanity check.

Also log the batch-mean histogram of $|s_\tau|$: if the bulk of
trajectories satisfy $|s_\tau| > \varepsilon$, V2/V3b can't possibly
differ from V1 (the spike region is empty).

\subsection*{(M2) Matched ablation runs (cost: $3\times$ Phase-1 compute)}

Fix all hyperparameters, seed, data, compute profile. Run three 120-step
BSPO runs with V1, V2a($\varepsilon=10^{-4}$), V3c respectively. Compare:
\begin{itemize}
\item Final \verb|val-core/<source>/acc/mean@1| for each of the 3
eval sources.
\item Training stability: $\mathrm{std}(\texttt{actor/pg\_loss})$ over
last $20$ steps, $\max(\texttt{actor/grad\_norm})$ over the whole run.
\item $\texttt{actor/bspo\_s\_tau\_abs\_mean}$ trajectory: expectation is
V2 pushes $|s_\tau|$ up more aggressively early; V1 lets $|s_\tau|$
rise organically from the magnitude path.
\end{itemize}

\subsection*{(M3) Synthetic backward-pass test (cost: minutes)}

Fabricate a batch where $s_\tau$ is tiny (say $10^{-5}$) and
$\min(s_\tau^+, s_\tau^-)$ is large (say $10^{-2}$). Run forward +
backward under each variant and compare $\partial\texttt{pg\_loss}/\partial\texttt{log\_prob}$
element-wise. V2a with $\varepsilon = 10^{-4}$ will produce gradients
$\sim 10^3\times$ larger than V1 on this batch. This establishes the
\emph{worst-case} injection under V2, but not the expected-case
contribution during real training.

% =============================================================================
\section{Recommendation}
\label{sec:reco}

\begin{enumerate}
\item \textbf{Default: V1 (\texttt{.detach()}).} Matches the mathematical
a.e.\ derivative of $\sign$, has no hyperparameter, introduces no
spurious gradient injection, and --- by
Theorem~\ref{thm:preannihilation} --- the on-policy factorisation
$\Delta^{(W)} = \sign\cdot m$ already makes the sign path contribute
zero at the fixed point. No optimisation benefit is available from
adding a surrogate.

\item \textbf{Switch to V2a only if V1 exhibits specific pathologies} ---
e.g.\ the policy getting stuck on the on-policy manifold and failing to
move. In the current Phase-1 ablation this pathology is
\emph{not observed}: $\texttt{actor/bspo\_s\_tau\_abs\_mean}$
rises monotonically from step $1$ onward, driven by the magnitude path.

\item \textbf{V3a is V1.} V3b and V3c are biased STE variants justified
only when a deliberate push off the on-policy manifold is part of the
algorithmic design; this is not the case here.

\item \textbf{If rigour is needed, use M1 first.} Log
$\|\text{grad}_{\text{sgn}}\|$ and $\|\text{grad}_{\text{mag}}\|$
separately. If the sign-path norm stays at zero under V1 and val-acc
curves are smooth, empirical confirmation of the theoretical
pre-annihilation argument is obtained.
\end{enumerate}

\textbf{One-sentence summary.} Since the magnitude factor
$\clip(\min(s_z^+, s_z^-), \delta_z)$ vanishes on the on-policy manifold
and the sign factor contributes only via this vanishing magnitude, the
sign path is pre-annihilated there regardless of how its gradient is
defined --- making V1's ``zero gradient through sign'' the
mathematically honest choice.

\end{document}
